% Declaração de Uso de Inteligência Artificial - UFAPE
% X marca a caixa; {} deixa em branco.
\documentclass[11pt,a4paper]{article}

% BLOCO DE DADOS: ALTERE APENAS OS CAMPOS ABAIXO
% Cabeçalho institucional, terceira linha
% Nome do curso, ex.: Bacharelado em Ciência da Computação
\newcommand{\FCursoHeader}{Bacharelado em Ciência da Computação}

% Identificação do manuscrito
% Título do trabalho, ex.: Uso de IA na revisão de textos acadêmicos
\newcommand{\FTitulo}{}
% Autores separados por vírgula, sem vírgula no fim; o primeiro é o correspondente
% Ex.: Maria da Silva, João Souza, Ana Lima
\newcommand{\FAutores}{}

% Ferramentas e propósitos
% Ferramentas separadas por vírgula, ex.: ChatGPT, Grammarly, Copilot
\newcommand{\FFerramentas}{}
% Revisão de linguagem, X marca
\newcommand{\FRev}{}
% Tradução, X marca
\newcommand{\FTra}{}
% Geração de texto, X marca
\newcommand{\FGer}{}
% Análise de dados, X marca
\newcommand{\FAna}{}
% Visualização de dados ou imagens, X marca
\newcommand{\FVis}{}
% Sugestões metodológicas, X marca
\newcommand{\FMet}{}
% Outros, X marca
\newcommand{\FOut}{}
% Detalhe dos outros, ex.: reescrita de código
\newcommand{\FOutDetalhe}{}

% Local e data, vão na última linha
% Cidade, ex.: Garanhuns
\newcommand{\FCidade}{Garanhuns}
% Dia, em número, ex.: 10
\newcommand{\FDia}{}
% Mês, por extenso, ex.: Janeiro
\newcommand{\FMes}{}
% Ano, em número, ex.: 2026
\newcommand{\FAno}{}

\usepackage[T1]{fontenc}
\usepackage{mathptmx}
\usepackage[brazilian]{babel}
\usepackage{graphicx}
\usepackage{setspace}
\usepackage{amssymb}
\usepackage{amsmath}
\usepackage{xcolor}
\usepackage{etoolbox}
\usepackage[normalem]{ulem}
\renewcommand{\ULthickness}{0.4pt}
\hyphenation{ma-tri-cu-la-do Con-clu-sao}
\usepackage[a4paper,left=3cm,right=3cm,top=1.6cm,bottom=1.6cm]{geometry}
\usepackage[bookmarksnumbered=true,psdextra,pdfencoding=auto]{hyperref}
\hypersetup{hidelinks}

\pagestyle{empty}
\setlength{\parindent}{0pt}
\setlength{\parskip}{0pt}
\setlength{\emergencystretch}{1em}

\newcommand{\gap}[2]{\vspace{#2mm}}

\newcommand{\caixa}[1]{%
  \raisebox{0.9pt}{$\vcenter{\hbox{%
    \setlength{\fboxsep}{1pt}\setlength{\fboxrule}{0.6pt}%
    \ifdefempty{#1}%
      {\fcolorbox{black}{white}{\phantom{{\large$\pmb{\checkmark}$}}}}%
      {\fcolorbox{black}{black!55}{\color{white}{\large$\pmb{\checkmark}$}}}%
  }}$}}

\newsavebox{\boxUF}

\newcommand{\UFaux}[1]{\uline{#1\hspace*{\fill}}}
\newcommand{\UF}[1]{\expandafter\UFaux\expandafter{#1}}

\newcommand{\restoUF}{}

\def\testespacoUF#1 #2\nilfim{\def\restoUF{#2}}

\newcommand{\campogrande}[2]{%
  \sbox{\boxUF}{\parbox{#1}{\hfuzz=100pt #2}}%
  \ifdim\ht\boxUF>1.5\baselineskip
    \noindent\parbox{#1}{#2}%
  \else
    \noindent\uline{\resizebox{#1}{!}{#2}}%
  \fi
  \par}

\newcommand{\campodinamico}[2]{%
  \expandafter\testespacoUF#2 \nilfim
  \ifx\restoUF\empty
    \sbox{\boxUF}{#1\hspace{0.15cm}#2}%
    \ifdim\wd\boxUF>\linewidth\relax
      \noindent#1\par
      \campogrande{\linewidth}{#2}%
    \else
      \noindent#1\hspace{0.15cm}\UF{#2}\par
    \fi
  \else
    \noindent#1\hspace{0.15cm}\UF{#2}\par
  \fi}

\newcommand{\flinha}[2]{\campodinamico{#1}{#2}}

\newcommand{\textoautores}{}

\newcommand{\acumulaaut}[1]{%
  \ifblank{#1}{}{%
    \ifdefempty{\textoautores}{}{\gappto\textoautores{; }}%
    \gappto\textoautores{#1}}}

\newcommand{\montaautores}{%
  \gdef\textoautores{}%
  \renewcommand*{\do}{\acumulaaut}%
  \expandafter\docsvlist\expandafter{\FAutores}}

\newcommand{\flinhaautores}{%
  \montaautores
  \campodinamico{Autor(es):}{\textoautores}}

\newcommand{\propitem}[3]{%
  \noindent\makebox[0.6cm][c]{\caixa{#1}}%
  \hspace{1mm}\mbox{\textbf{#2:} #3}\par}

\newcommand{\propitemlivre}[3]{%
  \noindent\makebox[0.6cm][c]{\caixa{#1}}\hspace{1mm}%
  \expandafter\testespacoUF#3 \nilfim
  \ifx\restoUF\empty
    \sbox{\boxUF}{\hspace*{0.7cm}\textbf{#2:}\hspace{0.15cm}#3}%
    \ifdim\wd\boxUF>\linewidth\relax
      \textbf{#2:}\par
      \campogrande{\linewidth}{#3}%
    \else
      \textbf{#2:}\hspace{0.15cm}\UF{#3}\par
    \fi
  \else
    \textbf{#2:}\hspace{0.15cm}\UF{#3}\par
  \fi}

\newcommand{\LAss}{7.2cm}

\newcommand{\assvazia}[1]{%
  \begin{minipage}[t]{\LAss}\centering
    \rule{\LAss}{0.4pt}\\
    \textbf{\small #1}
  \end{minipage}}

\newcommand{\fileira}[2]{%
  \noindent\assvazia{#1}\hfill\assvazia{#2}\par}

\newcommand{\sozinha}[1]{%
  \noindent\makebox[\textwidth][c]{\assvazia{#1}}\par}

\newcounter{idxsig}

\newcommand{\pendleg}{}

\newcommand{\linhasig}[1]{%
  \stepcounter{idxsig}%
  \ifodd\value{idxsig}%
    \ifnum\value{idxsig}=1 \gdef\pendleg{Assinatura do(a) autor(a) correspondente}%
    \else \gdef\pendleg{Assinatura do(a) coautor(a)}\fi
  \else
    \ifnum\value{idxsig}>2 \gap{ass-ass}{16.3}\fi
    \fileira{\pendleg}{Assinatura do(a) coautor(a)}%
  \fi}

\newcommand{\blocoassinaturas}{%
  \setcounter{idxsig}{0}%
  \renewcommand*{\do}{\linhasig}%
  \expandafter\docsvlist\expandafter{\FAutores}%
  \ifnum\value{idxsig}=0 \sozinha{Assinatura do(a) autor(a) correspondente}\fi
  \ifodd\value{idxsig}%
    \ifnum\value{idxsig}=1 \sozinha{Assinatura do(a) autor(a) correspondente}%
    \else \gap{ass-ass}{16.3}\sozinha{Assinatura do(a) coautor(a)}\fi
  \fi}

\begin{document}

{\setlength{\unitlength}{1mm}%
\begin{picture}(0,0)
  \put(-17.3,-10.4){\includegraphics[height=2.20cm]{figuras/brasao-ufape.pdf}}%
  \put(144,-10.4){\includegraphics[height=2.20cm]{figuras/marca-mec.pdf}}%
\end{picture}}%
\noindent\makebox[\textwidth][c]{%
  \begin{minipage}{14.2cm}\centering\bfseries\fontsize{11}{18.1}\selectfont
    Ministério da Educação\\
    Universidade Federal do Agreste de Pernambuco\\
    \ifdefempty{\FCursoHeader}{}{\FCursoHeader}%
  \end{minipage}%
}

\gap{cab-titulo}{3.5}
\noindent\makebox[\textwidth][c]{%
  \begin{minipage}{15cm}\centering\bfseries\fontsize{12}{15}\selectfont
    Declaração de Uso de Inteligência Artificial
  \end{minipage}%
}

\gap{titulo-1}{3.0}

\pdfbookmark[1]{Identificacao do manuscrito}{secId}
\flinha{Título do Trabalho:}{\FTitulo}

\gap{id-id}{1.8}
\flinhaautores

\gap{id-corpo}{1.2}
\begin{spacing}{1.22}
Nós, os autores do manuscrito acima identificado, declaramos que:
\end{spacing}

\gap{id-1}{3.0}
\pdfbookmark[1]{1 - Uso de Ferramentas de Inteligencia Artificial}{sec1}
\noindent\textbf{1- Uso de Ferramentas de Inteligência Artificial:}

\gap{1}{0.5}
\begin{spacing}{1.22}
Durante o desenvolvimento deste trabalho, as seguintes ferramentas de inteligência artificial foram utilizadas:
\end{spacing}

\gap{p1-fer}{1.2}
\flinha{Ferramenta(s):}{\FFerramentas}

\gap{fer-prop}{1.0}
\noindent\textbf{Propósito(s)} (marque todos os que se aplicam e forneça detalhes adicionais se necessário):

\gap{prop-itens}{0.5}
\propitem{\FRev}{Revisão de linguagem}{gramática, ortografia, clareza e estilo.}
\propitem{\FTra}{Tradução}{conversão do manuscrito para outro idioma.}
\propitem{\FGer}{Geração de texto}{rascunhos e sugestões revisados pelos autores.}
\propitem{\FAna}{Análise de dados}{processamento e interpretação de dados.}
\propitem{\FVis}{Visualização de dados ou imagens}{criação ou refinamento de figuras e gráficos.}
\propitem{\FMet}{Sugestões metodológicas}{ideias para aprimorar a pesquisa.}
\propitemlivre{\FOut}{Outros (especificar)}{\FOutDetalhe}

\gap{itens-2}{3.0}
\pdfbookmark[1]{2 - Autoria e Supervisao Humana}{sec2}
\noindent\textbf{2- Autoria e Supervisão Humana:}

\gap{2}{0.5}
\begin{spacing}{1.22}
A utilização de ferramentas de inteligência artificial foi exclusivamente para os propósitos acima declarados. Toda supervisão, validação e autoria intelectual do manuscrito são de responsabilidade dos autores humanos.
\end{spacing}

\gap{2-3}{3.0}
\pdfbookmark[1]{3 - Conformidade com Principios Eticos}{sec3}
\noindent\textbf{3- Conformidade com Princípios Éticos:}

\gap{3}{0.5}
\begin{spacing}{1.22}
O uso das ferramentas de IA seguiu práticas éticas e não incluiu atividades como fabricação de dados, plágio, manipulação de figuras ou qualquer prática que comprometa a integridade científica.
\end{spacing}

\gap{3-4}{3.0}
\pdfbookmark[1]{4 - Transparencia}{sec4}
\noindent\textbf{4- Transparência:}

\gap{4}{0.5}
\begin{spacing}{1.22}
Declaramos que as informações acima são completas e verdadeiras. Caso seja identificada qualquer inconsistência ou uso inadequado, aceitamos que o manuscrito seja submetido a investigação e, se necessário, rejeitado ou retratado.
\end{spacing}

\gap{texto-ass}{12.3}

\pdfbookmark[1]{Assinaturas}{secAss}
\noindent\begin{minipage}[t]{\textwidth}
\blocoassinaturas

\gap{ass-data}{11.3}
\noindent\makebox[\textwidth][c]{%
  \FCidade, \FDia\ de \FMes\ de \FAno.}
\end{minipage}

\end{document}
